% Options for packages loaded elsewhere
\PassOptionsToPackage{unicode}{hyperref}
\PassOptionsToPackage{hyphens}{url}
\documentclass[
  ignorenonframetext,
]{beamer}
\newif\ifbibliography
\usepackage{pgfpages}
\setbeamertemplate{caption}[numbered]
\setbeamertemplate{caption label separator}{: }
\setbeamercolor{caption name}{fg=normal text.fg}
\beamertemplatenavigationsymbolsempty
% remove section numbering
\setbeamertemplate{part page}{
  \centering
  \begin{beamercolorbox}[sep=16pt,center]{part title}
    \usebeamerfont{part title}\insertpart\par
  \end{beamercolorbox}
}
\setbeamertemplate{section page}{
  \centering
  \begin{beamercolorbox}[sep=12pt,center]{section title}
    \usebeamerfont{section title}\insertsection\par
  \end{beamercolorbox}
}
\setbeamertemplate{subsection page}{
  \centering
  \begin{beamercolorbox}[sep=8pt,center]{subsection title}
    \usebeamerfont{subsection title}\insertsubsection\par
  \end{beamercolorbox}
}
% Prevent slide breaks in the middle of a paragraph
\widowpenalties 1 10000
\raggedbottom
\AtBeginPart{
  \frame{\partpage}
}
\AtBeginSection{
  \ifbibliography
  \else
    \frame{\sectionpage}
  \fi
}
\AtBeginSubsection{
  \frame{\subsectionpage}
}
\usepackage{iftex}
\ifPDFTeX
  \usepackage[T1]{fontenc}
  \usepackage[utf8]{inputenc}
  \usepackage{textcomp} % provide euro and other symbols
\else % if luatex or xetex
  \usepackage{unicode-math} % this also loads fontspec
  \defaultfontfeatures{Scale=MatchLowercase}
  \defaultfontfeatures[\rmfamily]{Ligatures=TeX,Scale=1}
\fi
\usepackage{lmodern}
\ifPDFTeX\else
  % xetex/luatex font selection
\fi
% Use upquote if available, for straight quotes in verbatim environments
\IfFileExists{upquote.sty}{\usepackage{upquote}}{}
\IfFileExists{microtype.sty}{% use microtype if available
  \usepackage[]{microtype}
  \UseMicrotypeSet[protrusion]{basicmath} % disable protrusion for tt fonts
}{}
\makeatletter
\@ifundefined{KOMAClassName}{% if non-KOMA class
  \IfFileExists{parskip.sty}{%
    \usepackage{parskip}
  }{% else
    \setlength{\parindent}{0pt}
    \setlength{\parskip}{6pt plus 2pt minus 1pt}}
}{% if KOMA class
  \KOMAoptions{parskip=half}}
\makeatother
\usepackage{color}
\usepackage{fancyvrb}
\newcommand{\VerbBar}{|}
\newcommand{\VERB}{\Verb[commandchars=\\\{\}]}
\DefineVerbatimEnvironment{Highlighting}{Verbatim}{commandchars=\\\{\}}
% Add ',fontsize=\small' for more characters per line
\newenvironment{Shaded}{}{}
\newcommand{\AlertTok}[1]{\textcolor[rgb]{1.00,0.00,0.00}{\textbf{#1}}}
\newcommand{\AnnotationTok}[1]{\textcolor[rgb]{0.38,0.63,0.69}{\textbf{\textit{#1}}}}
\newcommand{\AttributeTok}[1]{\textcolor[rgb]{0.49,0.56,0.16}{#1}}
\newcommand{\BaseNTok}[1]{\textcolor[rgb]{0.25,0.63,0.44}{#1}}
\newcommand{\BuiltInTok}[1]{\textcolor[rgb]{0.00,0.50,0.00}{#1}}
\newcommand{\CharTok}[1]{\textcolor[rgb]{0.25,0.44,0.63}{#1}}
\newcommand{\CommentTok}[1]{\textcolor[rgb]{0.38,0.63,0.69}{\textit{#1}}}
\newcommand{\CommentVarTok}[1]{\textcolor[rgb]{0.38,0.63,0.69}{\textbf{\textit{#1}}}}
\newcommand{\ConstantTok}[1]{\textcolor[rgb]{0.53,0.00,0.00}{#1}}
\newcommand{\ControlFlowTok}[1]{\textcolor[rgb]{0.00,0.44,0.13}{\textbf{#1}}}
\newcommand{\DataTypeTok}[1]{\textcolor[rgb]{0.56,0.13,0.00}{#1}}
\newcommand{\DecValTok}[1]{\textcolor[rgb]{0.25,0.63,0.44}{#1}}
\newcommand{\DocumentationTok}[1]{\textcolor[rgb]{0.73,0.13,0.13}{\textit{#1}}}
\newcommand{\ErrorTok}[1]{\textcolor[rgb]{1.00,0.00,0.00}{\textbf{#1}}}
\newcommand{\ExtensionTok}[1]{#1}
\newcommand{\FloatTok}[1]{\textcolor[rgb]{0.25,0.63,0.44}{#1}}
\newcommand{\FunctionTok}[1]{\textcolor[rgb]{0.02,0.16,0.49}{#1}}
\newcommand{\ImportTok}[1]{\textcolor[rgb]{0.00,0.50,0.00}{\textbf{#1}}}
\newcommand{\InformationTok}[1]{\textcolor[rgb]{0.38,0.63,0.69}{\textbf{\textit{#1}}}}
\newcommand{\KeywordTok}[1]{\textcolor[rgb]{0.00,0.44,0.13}{\textbf{#1}}}
\newcommand{\NormalTok}[1]{#1}
\newcommand{\OperatorTok}[1]{\textcolor[rgb]{0.40,0.40,0.40}{#1}}
\newcommand{\OtherTok}[1]{\textcolor[rgb]{0.00,0.44,0.13}{#1}}
\newcommand{\PreprocessorTok}[1]{\textcolor[rgb]{0.74,0.48,0.00}{#1}}
\newcommand{\RegionMarkerTok}[1]{#1}
\newcommand{\SpecialCharTok}[1]{\textcolor[rgb]{0.25,0.44,0.63}{#1}}
\newcommand{\SpecialStringTok}[1]{\textcolor[rgb]{0.73,0.40,0.53}{#1}}
\newcommand{\StringTok}[1]{\textcolor[rgb]{0.25,0.44,0.63}{#1}}
\newcommand{\VariableTok}[1]{\textcolor[rgb]{0.10,0.09,0.49}{#1}}
\newcommand{\VerbatimStringTok}[1]{\textcolor[rgb]{0.25,0.44,0.63}{#1}}
\newcommand{\WarningTok}[1]{\textcolor[rgb]{0.38,0.63,0.69}{\textbf{\textit{#1}}}}
\setlength{\emergencystretch}{3em} % prevent overfull lines
\providecommand{\tightlist}{%
  \setlength{\itemsep}{0pt}\setlength{\parskip}{0pt}}
% Slide layout defaults for warning-clean scientific decks.
\usepackage{etoolbox}
\IfFileExists{xurl.sty}{\usepackage{xurl}}{}
\IfFileExists{seqsplit.sty}{\usepackage{seqsplit}}{\newcommand{\seqsplit}[1]{#1}}
\protected\def\breakseq#1{\seqsplit{#1}}
\protected\def\breaktt#1{\begingroup\ttfamily\seqsplit{#1}\endgroup}
\setlength{\emergencystretch}{6em}
\tolerance=5000
\hbadness=10000
\hfuzz=1pt
\setlength{\tabcolsep}{2pt}
\AtBeginEnvironment{longtable}{\tiny\renewcommand{\arraystretch}{0.86}\setlength{\tabcolsep}{1pt}}
\AtBeginEnvironment{tabular}{\tiny\renewcommand{\arraystretch}{0.86}\setlength{\tabcolsep}{1pt}}

% Natbib and cross-reference fallbacks — slides don't load natbib
% or cleveref, but combined-PDF manuscript prose may emit these
% commands. The fallback renders citations as a bracketed key list
% and cross-references as detokenized labels so slides don't fail on
% undefined control sequences or raw underscores. \providecommand is
% a no-op if packages load later.
\providecommand{\citep}[1]{[#1]}
\providecommand{\citet}[1]{#1}
\providecommand{\citealp}[1]{#1}
\providecommand{\citeauthor}[1]{#1}
\providecommand{\citeyear}[1]{#1}
\providecommand{\cref}[1]{\texttt{\detokenize{#1}}}
\providecommand{\Cref}[1]{\texttt{\detokenize{#1}}}

\newtheorem{proposition}{Proposition}
\newtheorem{hypothesis}{Hypothesis}
\usepackage{bookmark}
\IfFileExists{xurl.sty}{\usepackage{xurl}}{} % add URL line breaks if available
\urlstyle{same}
% fallback for those not using the hyperref driver hyperxmp:
\makeatletter
\@ifundefined{xmpquote}{\newcommand{\xmpquote}[1]{#1}}{}
\makeatother
\hypersetup{
  hidelinks,
  pdfcreator={LaTeX via pandoc}}

\author{\texorpdfstring{}{}}
\date{}

\begin{document}

\begin{frame}[fragile,allowframebreaks]{LLM-Based Assertion Extraction:
Prompt Design, Error Taxonomy, and Validation
\label{sec:extraction_pipeline}}
\protect\phantomsection\label{llm-based-assertion-extraction-prompt-design-error-taxonomy-and-validation}
\emph{This supplementary section documents the implementation specifics
of the LLM-based assertion extraction pipeline.}

\begin{block}{Relationship to Prior Approaches}
\protect\phantomsection\label{relationship-to-prior-approaches}
The closest prior effort is the systematic literature analysis of
Knight, Cordes, and Friedman \citep{knight2022fep}, which used human
annotators to manually code structural, visual, and mathematical
features of FEP and Active Inference publications. Their work operated
at the scale of hundreds of annotated papers and employed terms from the
Active Inference Institute's Active Inference Ontology for automated
text analysis. Our pipeline replaces the manual coding step with
LLM-based assertion extraction, enabling scalable processing of the full
corpus (\(N = 817\) papers) at the cost of exchanging human-verified
precision for machine-generated assessments that require post-hoc
validation. This trade-off is characteristic of the broader LLM-based
scientific extraction landscape: recent benchmarking confirms that even
state-of-the-art modular extraction architectures fall short of
production-level precision---particularly on tasks requiring exhaustive
retrieval and aggregation of multiple values from long
documents---validating our design choice to retain human review pathways
alongside automated extraction.

\begin{table}[htbp]
\centering
\caption{Comparison of annotation approaches: Knight et al.\ (2022) manual coding versus this work's automated LLM-based extraction pipeline. Key trade-offs between human-verified precision and machine-generated scalability are highlighted.}
\label{tab:annotation_comparison}
\begin{tabular}{lll}
\toprule
\textbf{Dimension} & \textbf{Knight et al.\ (2022)} & \textbf{This work} \\
\midrule
Scale & Hundreds of papers & 817 papers \\
Annotation & Manual (structural/visual/math features) & Automated (LLM hypothesis assessment) \\
Ontology & Active Inference Ontology terms & 8 standard hypotheses \\
Output & Annotated features + term frequencies & Nanopublications + knowledge graph \\
Reproducibility & Annotator-dependent & Deterministic (given model + seed) \\
Precision & High (human-verified) & Medium (requires validation) \\
\bottomrule
\end{tabular}
\end{table}

\begin{block}{Positioning in the LLM-Based Review Landscape}
\protect\phantomsection\label{positioning-in-the-llm-based-review-landscape}
Our pipeline operates within a rapidly maturing ecosystem of LLM-powered
literature analysis tools. Multi-agent architectures such as LitLLM
decompose the review process into specialized sub-agents (planner,
identifier, extractor, compiler), while ensemble approaches aggregate
outputs from multiple LLMs via weighted voting to improve reliability.
Our work differs from these tools in three respects: (1) we target
\emph{hypothesis-level evidence scoring} rather than inclusion/exclusion
screening; (2) we produce structured nanopublications rather than
narrative summaries; and (3) we are only analyzing abstracts for claims.
This deliberate trade-off enables corpus-scale processing (\(N = 817\))
but fundamentally misses fine-grained claims embedded in method sections
or discussion paragraphs. Full-text processing could improve extraction
recall, particularly for hypotheses with small evidence bases (H6
Clinical Utility, H7 Morphogenesis).
\end{block}
\end{block}

\begin{block}{The Eight Tracked Hypotheses}
\protect\phantomsection\label{the-eight-tracked-hypotheses}
Our analysis tracks the evolving evidence base for eight distinct claims
within the Active Inference literature, spanning theoretical
universality to applied clinical utility:

\begin{enumerate}
\tightlist
\item
  \textbf{H1: FEP Universality (Theoretical).} The Free Energy Principle
  applies universally to all self-organizing systems.
\item
  \textbf{H2: AIF Optimality (Computational).} Active Inference agents
  achieve optimal decision-making under uncertainty.
\item
  \textbf{H3: Markov Blanket Realism (Philosophical).} Markov blankets
  correspond to real physical boundaries.
\item
  \textbf{H4: Predictive Coding (Empirical).} Cortical hierarchies
  minimize prediction errors via predictive coding.
\item
  \textbf{H5: Scalability (Computational).} Active Inference scales to
  complex, high-dimensional environments.
\item
  \textbf{H6: Clinical Utility (Applied).} Active Inference provides
  clinically useful models of psychiatric conditions.
\item
  \textbf{H7: Morphogenesis (Biological).} The FEP explains
  morphogenetic and developmental processes.
\item
  \textbf{H8: Language AIF (Applied).} Active Inference provides a
  viable framework for language processing.
\end{enumerate}
\end{block}

\begin{block}{Prompt Engineering and Schema Design}
\protect\phantomsection\label{prompt-engineering-and-schema-design}
The structured prompt is designed to minimize parsing failures and
maximize assessment quality:

\begin{enumerate}
\item
  \textbf{Explicit JSON schema.} The prompt specifies the exact output
  schema---field names, allowed direction values, and the numeric
  confidence range---reducing the LLM's tendency to generate free-form
  text or ad hoc structures.
\item
  \textbf{Hypothesis definitions in-context.} All eight definitions are
  included verbatim, ensuring the LLM assesses relevance from the
  provided context rather than relying on parametric knowledge that may
  be stale.
\item
  \textbf{Reasoning field.} Each assessment includes a natural-language
  reasoning string, providing an audit trail for human reviewers and
  enabling systematic analysis of error patterns.
\item
  \textbf{Irrelevant filtering.} An explicit ``irrelevant'' direction
  allows the LLM to mark hypotheses that a paper does not address,
  avoiding forced spurious assessments.
\end{enumerate}

\begin{block}{Prompt Template}
\protect\phantomsection\label{prompt-template}
The extraction prompt follows a two-part structure (system + user):

\begin{Shaded}
\begin{Highlighting}[]
\NormalTok{SYSTEM: You are a scientific literature analyst specializing in the}
\NormalTok{Free Energy Principle and Active Inference. Assess the relevance of}
\NormalTok{the given paper to each hypothesis. Return a JSON array.}

\NormalTok{USER:}
\NormalTok{Paper: \{title\}}
\NormalTok{Abstract: \{abstract\}}

\NormalTok{Hypotheses:}
\NormalTok{H1: FEP Universality — \{description\}}
\NormalTok{H2: AIF Optimality — \{description\}}
\NormalTok{...}
\NormalTok{H8: Language AIF — \{description\}}

\NormalTok{For each hypothesis, return:}
\NormalTok{\{}
\NormalTok{  "hypothesis\_id": "H1",}
\NormalTok{  "direction": "supports|contradicts|neutral|irrelevant",}
\NormalTok{  "confidence": 0.0{-}1.0,}
\NormalTok{  "reasoning": "..."}
\NormalTok{\}}
\end{Highlighting}
\end{Shaded}

The extraction module (\breaktt{src/knowledge\_graph/llm\_extraction.py})
includes configurable retry logic with exponential backoff, JSON parsing
with handling of markdown code fences and extraneous text, confidence
clamping, and validation against the hypothesis ID set. The default
model is \texttt{gemma3:4b} on a local Ollama instance, configurable via
\texttt{-\/-llm-model} and \texttt{-\/-llm-url} flags.
\end{block}
\end{block}

\begin{block}{Failure Modes and Error Recovery}
\protect\phantomsection\label{failure-modes-and-error-recovery}
The primary failure modes are documented below.

\begin{block}{Over-Extraction Bias}
\protect\phantomsection\label{over-extraction-bias}
Preliminary experiments indicated \textasciitilde15--20\%
over-extraction. The current \(N = 817\) corpus was processed without a
validation set; error rates are not quantified for this run. This is the
most common error mode and produces false supporting evidence.
Over-extraction disproportionately affects broad-scope hypotheses (H1
FEP Universality, H2 AIF Optimality) where most papers in the corpus
contain relevant terminology without explicitly endorsing the claim.
Narrower hypotheses tied to specific domains (H7 Morphogenesis, H8
Language AIF) show lower over-extraction rates because their vocabulary
is more distinctive. This systematic bias inflates support counts for
broad hypotheses, and we caution against interpreting absolute scores
for H1 and H2 without accounting for this effect.
\end{block}

\begin{block}{Direction Misclassification}
\protect\phantomsection\label{direction-misclassification}
The LLM misclassifies a contradicting claim as supporting, or vice
versa. Rarer but more consequential, as it directly inverts the evidence
signal. Most common for papers that discuss limitations while ultimately
endorsing a hypothesis.
\end{block}

\begin{block}{Confidence Calibration Constraints}
\protect\phantomsection\label{confidence-calibration-constraints}
The model occasionally assigns high confidence to assessments where the
underlying evidence is ambiguous. Reliable confidence calibration
remains an open problem for zero-shot LLM applications, motivating the
multi-tiered validation protocols described below.
\end{block}

\begin{block}{Progressive JSON Parsing Recovery}
\protect\phantomsection\label{progressive-json-parsing-recovery}
To mitigate formatting inconsistencies, the module implements a
progressive parsing pipeline to recover malformed LLM outputs:

\begin{enumerate}
\tightlist
\item
  \textbf{Direct parse}: Attempt \texttt{json.loads()} on the raw
  response.
\item
  \textbf{Strip code fences}: Remove Markdown
  \texttt{\textasciigrave{}\textasciigrave{}\textasciigrave{}json\ ...\ \textasciigrave{}\textasciigrave{}\textasciigrave{}}
  wrappers and retry.
\item
  \textbf{Extract JSON array}: Scan for the first \texttt{{[}...{]}}
  substring in the response text.
\end{enumerate}

Papers that fail all parsing stages are logged and skipped; their count
is reported at pipeline completion.
\end{block}
\end{block}

\begin{block}{Validation Methodology}
\protect\phantomsection\label{validation-methodology}
To calibrate the pipeline in the absence of a human-annotated gold set,
we run a \textbf{rule-based reference-annotator agreement study} on a
stratified sample (\(n = 200\) assertions; hypothesis \(\times\) triage
direction \(\times\) year bin). Two \emph{deterministic
keyword-and-negation rule protocols}---not human annotators---label each
sampled abstract; their mutual agreement (Cohen's \(\kappa = 0.554\))
measures how stable the rule reference itself is, and the primary rule
protocol serves as an independent, fully reproducible reference against
which LLM pipeline triage is compared: precision 0.000, recall 0.000, F1
0.000 (positive class ``supports''). These figures are a reproducibility
floor, not an accuracy against ground truth: the rule reference and the
LLM diverge sharply (direction-agreement \(\kappa = -0.055\)), which is
itself the finding---absolute triage labels are unreliable and only
relative rankings should be trusted. Both protocols emit three separable
layers: source claim text, evidence supply (status and type), and
hypothesis triage direction. A human gold-standard annotation remains
future work (see the conclusion's ``Human spot-check coverage'').

\begin{table}[htbp]
\centering
\caption{Pipeline-versus-rule-reference agreement metrics (stratified sample). The reference labels are produced by deterministic keyword rules, not human annotators; values are a reproducibility floor, not accuracy against ground truth.}
\label{tab:validation_metrics}
\begin{tabular}{ll}
\toprule
\textbf{Metric} & \textbf{Value} \\
\midrule
Sample size & 200 \\
Inter-rule $\kappa$ (reference stability) & 0.554 \\
Reference--pipeline direction $\kappa$ & -0.055 \\
Pipeline precision vs.\ reference (supports) & 0.000 \\
Pipeline recall vs.\ reference (supports) & 0.000 \\
Pipeline F1 vs.\ reference (supports) & 0.000 \\
Over-extraction rate & 0.765 \\
\bottomrule
\end{tabular}
\end{table}

Error taxonomy rates (over-extraction, direction inversion, triage
mismatch) are reported in
\breaktt{output/reports/validation\_metrics.json}; the dominant mode is
over-extraction (0.765 of sampled rows), where the LLM assigns a
hypothesis label the keyword reference treats as irrelevant. Sensitivity
analysis across six citation-weight policies yields rank-stability
Spearman \(\rho = 0.762\) versus the default log-citation weighting,
with 18 hypothesis rank-position changes across all alternative
policies.
\end{block}

\begin{block}{From Assertions to Nanopublications}
\protect\phantomsection\label{from-assertions-to-nanopublications}
Each assertion is wrapped in a \textbf{nanopublication}
\citep{groth2010anatomy, kuhn2016decentralized} that carries a
structured provenance block---paper ID, source passage, model ID
(\texttt{gemma3:4b}), prompt version (\texttt{v2.0.4}), processing date,
pipeline version, and run ID. Every assertion in the corpus is produced
by the three-layer extractor and carries this block; the aggregate
model, prompt-version, and processing-date coverage is reported in
\breaktt{output/reports/extraction\_provenance\_summary.json}.

Nanopublications are persisted \textbf{incrementally} during extraction.
Every 50 papers (configurable via \breaktt{-\/-checkpoint-interval}), the
pipeline atomically appends newly extracted nanopublications to
\breaktt{nanopublications.jsonl} using a temporary-file-plus-rename
strategy that prevents corruption on interruption. Deduplication
operates on the composite key \((paper\_id, hypothesis\_id)\): when a
paper is re-processed with an improved model, the newer assertion
overwrites the stale entry. This merge-on-add design enables iterative
model refinement without costly full-corpus re-extraction.

After extraction completes, the full nanopublication set is additionally
serialized to \textbf{RDF/TriG} format per the nanopublication standard,
producing four named graphs per nanopublication (Head, Assertion,
Provenance, Publication Info). The TriG output is suitable for
publication to the decentralized nanopublication network and archival on
data repositories such as Zenodo. The complete RDF schema is specified
in the \hyperref[sec:methods_kg]{knowledge graph methodology} and
Appendix \ref{sec:appendix_rdf}.
\end{block}
\end{frame}

\end{document}
